\scanpage{115}
% archive foliation: 114
ouvert \inserted{$V$} de $U^{\mathrm{an}}$, sont des restrictions à $V$ de fonctions \inserted{\unsure{multiformes}} uniformisantes sur $X^{\mathrm{an}}$ (et on \unsure{trouve} \inserted{ainsi} loc. exactement toutes les fonctions uniformisantes).

\uline{Corollaire} Si $X$ n'est pas complète, \ldots le $=$ énoncé avec la seule \unsure{Diff.} que $\mu(t)$ est seulement une section \uline{analytique} de $\mathcal O_U$. \deleted{\unsure{Pouvons nous la}} \inserted{Est-elle} algébrique ? Premiers cas classiques : $X=\mathbb P^1_{\mathbb C}-\text{trois pts}$, ou $X=\text{courbe elliptique}$ moins \uline{un point}.

\uline{Autre question} Quand $X$ (supposons compact disons) varie avec un paramètre \uline{algébrique} parcourant une variété algébrique $S$, i.e. \ldots $X\longrightarrow S$, \unsure{ça donne} comme $\mu$ des fonctions analytiques sur un ouvert de $X$, algébriques sur les fibres. Sont-elles algébriques ? \unsure{Celles-ci, ont elles|Ceci a-t-il} une signification purement algébrique, indépendante de la top. du corps de base $\mathbb C$ -- donc ayant un sens dès que $S$ est un schéma de car nulle ?).
